\textbf{Phần 1. Trắc nghiệm lựa chọn}
\setcounterpageref{ex}{0}

\begin{ex} %[3D8N45]
Số gồm $4$ nghìn, $5$ trăm, $2$ chục và $7$ đơn vị viết là
\choice
{\True $4527$}
{$4257$}
{$5427$}
{$4572$}
\loigiai{
Số đó là $4527$.\\
\textbf{Đáp án của câu hỏi là A}
}
\end{ex}

\begin{ex} %[3D8N45]
Số $10\,000$ đọc là
\choice
{Một nghìn}
{\True Mười nghìn (hay một vạn)}
{Một trăm nghìn}
{Mười triệu}
\loigiai{
Số $10\,000$ đọc là Mười nghìn.\\
\textbf{Đáp án của câu hỏi là B}
}
\end{ex}

\begin{ex} %[3D8H45]
Trong số $7\,845$, chữ số 7 có giá trị là
\choice
{$700$}
{$70$}
{\True $7000$}
{$7$}
\loigiai{
Chữ số 7 thuộc hàng nghìn nên có giá trị là $7000$.\\
\textbf{Đáp án của câu hỏi là C}
}
\end{ex}

\begin{ex} %[3D8H45]
Số liền sau của số $9\,999$ là
\choice
{$9998$}
{\True $10000$}
{$1000$}
{$9990$}
\loigiai{
$9999 + 1 = 10\,000$.\\
\textbf{Đáp án của câu hỏi là B}
}
\end{ex}

\begin{ex} %[3D8H45]
Số liền trước của số $5\,000$ là
\choice
{$5001$}
{$4990$}
{$4900$}
{\True $4999$}
\loigiai{
$5000 - 1 = 4999$.\\
\textbf{Đáp án của câu hỏi là D}
}
\end{ex}

\begin{ex} %[3D8H45]
Số \textit{Ba nghìn không trăm linh sáu} viết là
\choice
{$3600$}
{\True $3006$}
{$3060$}
{$3606$}
\loigiai{
Số đó viết là $3006$.\\
\textbf{Đáp án của câu hỏi là B}
}
\end{ex}

\begin{ex} %[3D8V45]
Số gồm $6$ nghìn và $8$ đơn vị viết là
\choice
{\True $6008$}
{$6800$}
{$6080$}
{$6808$}
\loigiai{
Số đó là $6008$.\\
\textbf{Đáp án của câu hỏi là A}
}
\end{ex}

\begin{ex} %[3D8V45]
Viết số $4\,325$ thành tổng các hàng:
\choice
{$400 + 300 + 20 + 5$}
{$4000 + 30 + 20 + 5$}
{\True $4000 + 300 + 20 + 5$}
{$4000 + 320 + 5$}
\loigiai{
$4325 = 4000 + 300 + 20 + 5$.\\
\textbf{Đáp án của câu hỏi là C}
}
\end{ex}

\begin{ex} %[3D8V45]
Từ 4 chữ số $2, 0, 5, 8$ lập được số lớn nhất có 4 chữ số khác nhau là
\choice
{$8502$}
{$8250$}
{$5820$}
{\True $8520$}
\loigiai{
Số lớn nhất là $8520$.\\
\textbf{Đáp án của câu hỏi là D}
}
\end{ex}

\begin{ex} %[3D8V45]
Từ 4 chữ số $1, 0, 4, 7$ lập được số bé nhất có 4 chữ số khác nhau là
\choice
{$1407$}
{$1074$}
{\True $1047$}
{$0147$}
\loigiai{
Số bé nhất là $1047$.\\
\textbf{Đáp án của câu hỏi là C}
}
\end{ex}

\begin{ex} %[3D8V45]
Dãy số tròn nghìn từ $1\,000$ đến $9\,000$ có tất cả bao nhiêu số?
\choice
{\True $9$ số}
{$8$ số}
{$10$ số}
{$7$ số}
\loigiai{
Các số là $1000, 2000, \dots, 9000$ (có 9 số).\\
\textbf{Đáp án của câu hỏi là A}
}
\end{ex}

\begin{ex} %[3D8C45]
Tìm số có 4 chữ số biết chữ số hàng nghìn là 3, chữ số hàng trăm gấp 2 lần chữ số hàng nghìn, chữ số hàng chục và hàng đơn vị đều là 0.
\choice
{$3300$}
{$3800$}
{$3900$}
{\True $3600$}
\loigiai{
Hàng nghìn: 3. Hàng trăm: $3 \times 2 = 6$. Số đó là $3600$.\\
\textbf{Đáp án của câu hỏi là D}
}
\end{ex}

\textbf{Phần 2. Trắc nghiệm đúng sai}
\setcounterpageref{ex}{0}

\begin{ex} %[3D8H45]
Xét tính đúng sai của các nhận định về số có 4 chữ số:
\choiceTF
{\True $1000$ là số bé nhất có 4 chữ số.}
{\True $9999$ là số lớn nhất có 4 chữ số.}
{Số liền trước của $10000$ là $9990$.}
{\True Số liền trước của $10000$ là $9999$.}
\loigiai{
\begin{itemize}
    \item $1000$ bé nhất có 4 chữ số: Đúng.
    \item $9999$ lớn nhất có 4 chữ số: Đúng.
    \item $10000 - 1 = 9999 \neq 9990$. Ý 3 Sai, Ý 4 Đúng.
\end{itemize}
\begin{itemize}
    \item a) \textbf{Đúng}.
    \item b) \textbf{Đúng}.
    \item c) \textbf{Sai}.
    \item d) \textbf{Đúng}.
\end{itemize}
}
\end{ex}

\begin{ex} %[3D8V45]
Cho số $6\,504$. Xét tính đúng sai:
\choiceTF
{\True Chữ số 6 có giá trị là $6000$.}
{\True Chữ số 5 có giá trị là $500$.}
{Chữ số hàng chục là chữ số 4.}
{\True Chữ số hàng chục là chữ số 0, chữ số hàng đơn vị là 4.}
\loigiai{
\begin{itemize}
    \item Giá trị chữ số 6: $6000$. Đúng.
    \item Giá trị chữ số 5: $500$. Đúng.
    \item Hàng chục là 0 (không phải 4). Ý 3 Sai, Ý 4 Đúng.
\end{itemize}
\begin{itemize}
    \item a) \textbf{Đúng}.
    \item b) \textbf{Đúng}.
    \item c) \textbf{Sai}.
    \item d) \textbf{Đúng}.
\end{itemize}
}
\end{ex}

\textbf{Phần 3. Trả lời ngắn}
\setcounterpageref{ex}{0}

\begin{ex} %[3D8H45]
Viết số liền sau của số $4\,599$.
\shortans{4600}
\loigiai{
\textbf{Đáp số:} $4600$\\
$4599 + 1 = 4600$.
}
\end{ex}

\begin{ex} %[3D8H45]
Viết số gồm $7$ nghìn, $8$ trăm và $5$ đơn vị.
\shortans{7805}
\loigiai{
\textbf{Đáp số:} $7805$\\
$7805$.
}
\end{ex}

\begin{ex} %[3D8V45]
Tìm số lớn nhất có 4 chữ số khác nhau lập từ các chữ số $9, 3, 5, 1$.
\shortans{9531}
\loigiai{
\textbf{Đáp số:} $9531$\\
$9531$.
}
\end{ex}

\begin{ex} %[3D8V45]
Tìm số bé nhất có 4 chữ số khác nhau lập từ các chữ số $4, 0, 7, 2$.
\shortans{2047}
\loigiai{
\textbf{Đáp số:} $2047$\\
$2047$.
}
\end{ex}

\begin{ex} %[3D8V45]
Tính tổng giá trị các chữ số của số $3\,421$.
\shortans{10}
\loigiai{
\textbf{Đáp số:} $10$\\
$3 + 4 + 2 + 1 = 10$.
}
\end{ex}

\begin{ex} %[3D8C45]
Tìm số tròn trăm lớn nhất có 4 chữ số mà tổng các chữ số của nó bằng 12 (biết chữ số hàng nghìn là 7).
\shortans{7500}
\loigiai{
\textbf{Đáp số:} $7500$\\
Số tròn trăm có dạng $\overline{7a00}$. Ta có $7 + a + 0 + 0 = 12 \Rightarrow a = 5$. Số đó là $7500$.
}
\end{ex}
